\documentclass[12pt,a4paper]{article}
\usepackage[utf8]{inputenc}
\usepackage{amsmath, amssymb}
\usepackage{enumitem}
\usepackage[margin=2cm]{geometry}

\begin{document}

\begin{center}
    \textbf{\Large Latihan Soal PABP Kelas X} \\
    \textbf{Persiapan Asesmen Sumatif Akhir Tahun (ASAT)}
\end{center}
\vspace{0.5cm}

\begin{enumerate}[leftmargin=*]

% ==========================================
% MATERI: AKHLAK MADZMUMAH & MAHMUDAH 
% (Ghadab, Syaja'ah, Mujahadatun Nafs)
% ==========================================

\item Secara terminologi akademik, pengertian yang paling tepat mengenai sifat temperamental (\textit{ghadhab}) adalah \dots
    \begin{enumerate}[label=\alph*.]
        \item Sikap tenang dalam menghadapi segala bentuk intimidasi dan provokasi dari pihak luar.
        \item Gejolak emosional negatif di dalam hati yang menggerakkan darah untuk meluap, menuntut kepuasan agresi, serta melumpuhkan kontrol akal sehat.
        \item Kekuatan hati untuk melakukan perlawanan fisik demi menegakkan keadilan di lingkungan sosial.
        \item Rasa tidak suka yang disembunyikan di dalam dada karena ketidakmampuan membalas perbuatan buruk orang lain.
        \item Perilaku berani yang tidak terukur dalam merespons ketidakadilan secara langsung dan spontan.
    \end{enumerate}

\item Dalam konsep etika Islam, sikap berani membela kebenaran (\textit{syaja'ah}) didefinisikan sebagai titik keseimbangan (\textit{al-wasath}) yang ideal di antara dua kutub ekstrem yang tercela, yaitu \dots
    \begin{enumerate}[label=\alph*.]
        \item \textit{Riya'} (pamer) dan \textit{Sum'ah} (ingin didengar)
        \item \textit{Kibar} (sombong) dan \textit{Ujub} (merasa bangga diri)
        \item \textit{Jubn} (pengecut) dan \textit{Tahawwur} (nekat membabi buta)
        \item \textit{Ghadhab} (pemarah) dan \textit{Fahisyah} (keji)
        \item \textit{Ananiah} (egois) dan \textit{Takabur} (angkuh)
    \end{enumerate}

\item Adanya penyakit hati seperti kesombongan (\textit{kibar}), merasa diri paling benar (\textit{ujub}), egoisme yang tinggi, serta kurangnya latihan pengendalian diri merupakan faktor pencetus sifat \textit{ghadhab} yang bersumber dari dimensi \dots
    \begin{enumerate}[label=\alph*.]
        \item Internal (Psikologis)
        \item Eksternal (Lingkungan)
        \item Struktural (Sosial)
        \item Kultural (Budaya)
        \item Komunal (Kelompok)
    \end{enumerate}

\item Perjuangan bersungguh-sungguh melawan kecenderungan hawa nafsu internal yang buruk (\textit{nafsul ammarah bis-su'}) agar tunduk pada kendali akal sehat dan bimbingan syariat disebut \dots
    \begin{enumerate}[label=\alph*.]
        \item \textit{Syaja'ah}
        \item \textit{Mujahadatun Nafs}
        \item \textit{Hifzhu al-Din}
        \item \textit{Mahabbatullah}
        \item \textit{Tawakkal}
    \end{enumerate}

\item Salah satu fondasi spiritual yang paling utama untuk membentuk karakter \textit{syaja'ah} yang murni di dalam diri seorang mukmin adalah \dots
    \begin{enumerate}[label=\alph*.]
        \item Kekuatan fisik dan ketahanan tubuh yang terlatih sejak dini
        \item Keyakinan tauhid yang murni bahwa hidup, mati, dan rezeki berada sepenuhnya di bawah ketetapan Allah SWT
        \item Dukungan penuh dari kelompok sosial dan lingkungan masyarakat sekitar
        \item Orientasi material duniawi untuk meningkatkan status sosial pribadi
        \item Keinginan kuat untuk mendapatkan pujian atas keberanian yang ditunjukkan
    \end{enumerate}

\item Dari perspektif sosial, hikmah dan manfaat nyata dari tegaknya pilar kesalehan karakter \textit{syaja'ah} di tengah masyarakat adalah \dots
    \begin{enumerate}[label=\alph*.]
        \item Terwujudnya stabilitas keamanan hukum dan terlindunginya kaum lemah dari kezaliman struktural
        \item Lahirnya kelompok-kelompok militer yang berkuasa penuh atas tatanan sosial
        \item Hilangnya perbedaan pendapat di antara kelompok-kelompok keagamaan
        \item Terkumpulnya harta rampasan yang banyak untuk kemaslahatan ekonomi sepihak
        \item Meningkatnya populasi penduduk karena jaminan keamanan fisik yang mutlak
    \end{enumerate}

\item Keberanian sejati tidak lahir dari kekuatan fisik semata. Selain tauhid yang murni, ketabahan dalam menghadapi ujian berat dibentuk oleh hati yang bersih dari ketergantungan kepada makhluk, atau disebut juga dengan \dots
    \begin{enumerate}[label=\alph*.]
        \item \textit{Sakinah} (Ketenangan jiwa)
        \item \textit{Kibar} (Kesombongan)
        \item \textit{Fahisyah} (Kekejian)
        \item \textit{Ghadhab} (Kemarahan)
        \item \textit{Tamanni} (Anggan-angan kosong)
    \end{enumerate}

\item Melakukan jeda berpikir sebelum merespons stimulus eksternal agar tidak terjebak dalam penyesalan merupakan langkah konkret dalam pembiasaan kontrol diri pada aspek \dots
    \begin{enumerate}[label=\alph*.]
        \item Melatih disiplin ibadah puasa
        \item Muhasabah rutin akhir hari
        \item Manajemen emosi aktual
        \item Penapisan informasi digital
        \item Meditasi spiritual formal
    \end{enumerate}

\item Mengondisikan diri untuk menahan keinginan biologis yang sah (makan dan minum) bertujuan untuk mengasah ketajaman spiritual dan mengendalikan hawa nafsu secara terstruktur. Metode pembiasaan kontrol diri ini dilakukan melalui \dots
    \begin{enumerate}[label=\alph*.]
        \item Berwudhu dengan air mengalir
        \item Membaca kalimat \textit{ta'awwudz}
        \item Melatih disiplin lewat ibadah puasa
        \item Mengubah posisi tubuh saat berinteraksi
        \item Melakukan muhasabah rutin terhadap niat
    \end{enumerate}

\item Bentuk implementasi nyata dari sikap \textit{syaja'ah} ketika berhadapan dengan otoritas kepemimpinan yang bertindak menyimpang adalah \dots
    \begin{enumerate}[label=\alph*.]
        \item Menggalang massa untuk melakukan tindakan anarkis di fasilitas publik
        \item Berani menyampaikan kebenaran secara santun (\textit{kalimatu haqqin 'inda sulthanin ja'ir})
        \item Diam dan memendam kebenaran tersebut demi menjaga stabilitas pribadi
        \item Menyebarkan ujaran kebencian tanpa dasar yang valid melalui media sosial
        \item Memutus hubungan diplomatik secara sepihak tanpa adanya jalur dialog
    \end{enumerate}

% ==========================================
% MATERI: AL-KULLIYATU AL-KHAMSAH 
% ==========================================

\item Kajian Ushul Fikih dan filosofi hukum Islam mengenal istilah \textit{Al-Kulliyatu al-Khamsah}. Istilah universal tersebut memiliki arti \dots
    \begin{enumerate}[label=\alph*.]
        \item Lima prinsip ibadah pokok yang wajib dijalankan oleh setiap umat Islam
        \item Lima metode dakwah kultural yang digunakan dalam penyebaran Islam
        \item Lima kepentingan mutlak kemanusiaan yang wajib dijaga, dilindungi, dan dipelihara
        \item Lima tahapan pembentukan karakter mulia dalam diri seorang mukmin
        \item Lima larangan hukum pidana yang memiliki sanksi eksekusi berat di dunia
    \end{enumerate}

\item Membentengi diri dari pemikiran menyimpang (seperti sekularisme radikal dan liberalisme agama), menjauhi kesyirikan, serta menolak penistaan terhadap syariat merupakan metode pemeliharaan \textit{hifzhu al-din} dari \dots
    \begin{enumerate}[label=\alph*.]
        \item Aspek Konstruktif (\textit{Min janib al-wujud})
        \item Aspek Defensif (\textit{Min janib al-'adam})
        \item Aspek Preventif Kontemporer
        \item Aspek Regulasi Struktural
        \item Aspek Edukasi Komparatif
    \end{enumerate}

\item Menegakkan rukun Islam secara konsisten, memperdalam kajian akidah yang shahih, serta menyebarkan syiar dakwah secara damai merupakan upaya menjaga agama (\textit{hifzhu al-din}) melalui \dots
    \begin{enumerate}[label=\alph*.]
        \item Aspek Defensif (\textit{Min janib al-'adam})
        \item Aspek Struktural Kontekstual
        \item Aspek Konstruktif (\textit{Min janib al-wujud})
        \item Aspek Sanktioner Formal
        \item Aspek Kultural Komunal
    \end{enumerate}

\item Instrumen normatif dalam syariat Islam yang memberlakukan sanksi balasan setimpal (\textit{qishash}) dan denda materi (\textit{diyat}) dirancang secara tegas untuk merealisasikan prinsip universal \dots
    \begin{enumerate}[label=\alph*.]
        \item Menjaga Harta (\textit{Hifzhu al-Maal})
        \item Menjaga Akal (\textit{Hifzhu al-Aql})
        \item Menjaga Jiwa (\textit{Hifzhu al-Nafs})
        \item Menjaga Keturunan (\textit{Hifzhu al-Nasl})
        \item Menjaga Kehormatan (\textit{Hifzhu al-Irdh})
    \end{enumerate}

\item Upaya menjaga kesehatan tubuh melalui konsumsi makanan bergizi, pemenuhan layanan medis yang layak, dan pemeliharaan keselamatan lingkungan kerja termasuk dalam kategori \dots
    \begin{enumerate}[label=\alph*.]
        \item Kewajiban individu dalam menjaga harta pribadi
        \item Cara menjaga jiwa (\textit{hifzhu al-nafs}) demi kelangsungan hidup terhormat
        \item Penerapan rukun Islam pada dimensi kesehatan jasmani
        \item Strategi budaya untuk meningkatkan taraf ekonomi umat
        \item Komponen sekunder yang tidak memengaruhi orisinalitas iman
    \end{enumerate}

\item Landasan utama penetapan hukum kewajiban menuntut ilmu (\textit{thalabul 'ilmi}) secara kontinu dalam Islam ditujukan untuk mengimplementasikan \dots
    \begin{enumerate}[label=\alph*.]
        \item Pengembangan kapasitas intelektual untuk memelihara akal (\textit{hifzhu al-aql})
        \item Peningkatan keahlian dagang guna mencari harta (\textit{hifzhu al-maal})
        \item Penguatan status kepemimpinan politik di dalam lingkungan sosial
        \item Perlindungan keturunan dari kebodohan doktrin pergaulan bebas
        \item Pertahanan fisik dari agresi bersenjata luar
    \end{enumerate}

\item Pengharaman mutlak atas konsumsi \textit{khamr}, narkotika, psikotropika, dan zat adiktif lainnya, disertai penetapan sanksi hukum \textit{hadd}, merupakan langkah defensif syariat untuk melindungi \dots
    \begin{enumerate}[label=\alph*.]
        \item Eksistensi keturunan dan kesucian garis nasab keluarga
        \item Basis pembebanan syariat (\textit{manathat taklif}) berupa fungsi akal sehat
        \item Hak kepemilikan aset privat dan publik dari eksploitasi zalim
        \item Ketenangan jiwa spiritual dalam menyerap keindahan alam semesta
        \item Struktur tatanan hukum pidana dari pengaruh intervensi asing
    \end{enumerate}

\item Mensyariatkan institusi pernikahan yang sah dan suci sebagai gerbang legal-spiritual interaksi biologis merupakan bentuk perlindungan konkret terhadap \dots
    \begin{enumerate}[label=\alph*.]
        \item Keberadaan harta benda peninggalan orang tua
        \item Kelestarian eksistensi umat manusia dan kesucian nasab (\textit{hifzhu al-nasl})
        \item Kebebasan berpikir individual tanpa batas moralitas
        \item Stabilitas fisik manusia dari paparan penyakit medis menular
        \item Kehormatan institusi formal kedinasan pemerintahan
    \end{enumerate}

\item Kewajiban orang tua memberikan nafkah yang halal dan pendidikan akhlakul karimah yang komprehensif bagi anak cucunya merupakan aplikasi dari \dots
    \begin{enumerate}[label=\alph*.]
        \item Cara memelihara keturunan (\textit{hifzhu al-nasl}) agar berkarakter mulia
        \item Metode eksplorasi sumber daya ekonomi secara individual-kapitalis
        \item Cara menjaga rahasia internal urusan rumah tangga dari publik
        \item Strategi politik keluarga untuk memenangi persaingan birokrasi
        \item Upaya defensif melindungi aset finansial dari inflasi global
    \end{enumerate}

\item Pengharaman praktik pencurian, korupsi, manipulasi data keuangan, perjudian, serta transaksi berbasis riba bertujuan untuk menegakkan prinsip universal \dots
    \begin{enumerate}[label=\alph*.]
        \item Menjaga Jiwa (\textit{Hifzhu al-Nafs})
        \item Menjaga Akal (\textit{Hifzhu al-Aql})
        \item Menjaga Agama (\textit{Hifzhu al-Din})
        \item Menjaga Keturunan (\textit{Hifzhu al-Nasl})
        \item Menjaga Harta (\textit{Hifzhu al-Maal})
    \end{enumerate}

% ==========================================
% MATERI: DAKWAH WALI SONGO 
% ==========================================

\item Karakteristik utama dari proses Islamisasi di Nusantara yang diarsiteki oleh Dewan Wali Songo di tanah Jawa adalah keberhasilan penerapan metode \dots
    \begin{enumerate}[label=\alph*.]
        \item Dakwah struktural melalui jalur konfrontasi militer dan politik radikal
        \item Dakwah kultural (\textit{bil-hikmah} dan \textit{mau'izhatul hasanah}) berbasis pribumisasi Islam
        \item Dakwah eksklusif yang memisahkan diri sepenuhnya dari tradisi masyarakat lokal
        \item Dakwah materialistis dengan memberikan modal keuangan bersyarat kepada rakyat
        \item Dakwah doktriner kaku tanpa mempertimbangkan kondisi psikososial objek dakwah
    \end{enumerate}

\item Strategi pribumisasi Islam yang diterapkan oleh Wali Songo diwujudkan melalui cara \dots
    \begin{enumerate}[label=\alph*.]
        \item Menghapus total seluruh adat istiadat lama tanpa menyisakan unsur budaya lokal
        \item Mengintegrasikan nilai-nilai tauhid ke dalam kebudayaan lokal tanpa merusak struktur tradisi asli
        \item Mengadopsi sistem kasta sosial untuk mempermudah pemetaan target dakwah
        \item Membatasi interaksi sosial dengan penduduk lokal yang belum masuk Islam
        \item Mewajibkan penggunaan bahasa asing dalam seluruh upacara adat kemasyarakatan
    \end{enumerate}

\item Tokoh pelopor Wali Songo yang menggunakan pendekatan ekonomi melalui jalur perdagangan internasional, pembaruan sistem pertanian, serta membuka layanan pengobatan gratis bagi kasta Sudra adalah \dots
    \begin{enumerate}[label=\alph*.]
        \item Sunan Ampel
        \item Sunan Kalijaga
        \item Sunan Gresik (Maulana Malik Ibrahim)
        \item Sunan Kudus
        \item Sunan Giri
    \end{enumerate}

\item Sunan Ampel (Raden Rahmat) merupakan perancang \textit{blueprint} kaderisasi dai di tanah Jawa yang mendirikan Pesantren Ampeldenta. Beliau sangat terkenal dengan doktrin moral \textbf{"Moh Limo"}. Arti dari larangan \textit{Moh Madat} dan \textit{Moh Madon} berturut-turut adalah \dots
    \begin{enumerate}[label=\alph*.]
        \item Tidak berjudi dan tidak mabuk
        \item Tidak mencuri dan tidak berzina
        \item Tidak mengonsumsi narkoba dan tidak berzina
        \item Tidak mabuk dan tidak mencuri
        \item Tidak berjudi dan tidak mengonsumsi narkoba
    \end{enumerate}

\item Pendekatan dakwah yang bertumpu pada filantropi Islam, pemberdayaan ekonomi, jaminan sosial bagi kaum miskin dan anak yatim, serta penggagas media tembang \textit{Pangkur} merupakan karakteristik dari \dots
    \begin{enumerate}[label=\alph*.]
        \item Sunan Gunung Jati
        \item Sunan Drajat (Raden Qasim)
        \item Sunan Giri
        \item Sunan Kudus
        \item Sunan Ampel
    \end{enumerate}

\item Bentuk toleransi kultural yang sangat tinggi ditunjukkan oleh Sunan Kudus (Ja'far Shadiq) dalam merangkul psikologi budaya penganut Hindu-Buddha di wilayah dakwahnya, yaitu dengan cara \dots
    \begin{enumerate}[label=\alph*.]
        \item Mengubah nama kota menjadi wilayah bernuansa asing secara paksa
        \item Melarang penyembelihan sapi dan mendesain arsitektur menara masjid menyerupai candi
        \item Mewajibkan seluruh masyarakat lokal memakai pakaian khas wilayah pusat Islam
        \item Menolak sistem hukum fikih formal dalam menyelesaikan konflik sosial masyarakat
        \item Membongkar candi-candi lama untuk dijadikan bahan baku utama fondasi masjid
    \end{enumerate}

\item Wali Songo yang mengembangkan jalur dakwah berbasis institusi pendidikan pesantren Giri Kedaton dan kekuatan politik struktural, serta menciptakan permainan anak Islami seperti \textit{Jelungan} dan tembang \textit{Asmarandana} adalah \dots
    \begin{enumerate}[label=\alph*.]
        \item Sunan Kalijaga
        \item Sunan Gunung Jati
        \item Sunan Giri (Raden Paku)
        \item Sunan Gresik
        \item Sunan Drajat
    \end{enumerate}

\item Modifikasi mendalam terhadap kesenian wayang kulit, instrumen gamelan, dan seni ukir agar bermuatan falsafah tauhid merupakan strategi dakwah spesifik yang dilakukan oleh budayawan agung Wali Songo, yaitu \dots
    \begin{enumerate}[label=\alph*.]
        \item Sunan Kalijaga (Raden Mas Said)
        \item Sunan Kudus
        \item Sunan Ampel
        \item Sunan Gresik
        \item Sunan Gunung Jati
    \end{enumerate}

\item Sunan Kalijaga merupakan pencipta karya sastra tembang mistikus-edukatif legendaris di tanah Jawa yang berjudul \dots
    \begin{enumerate}[label=\alph*.]
        \item \textit{Pangkur} dan \textit{Asmarandana}
        \item \textit{Ilir-Ilir} dan \textit{Gundul-Gundul Pacul}
        \item \textit{Jelungan} dan \textit{Asmarandana}
        \item \textit{Moh Limo} dan \textit{Pangkur}
        \item \textit{Asmarandana} dan \textit{Ilir-Ilir}
    \end{enumerate}

\item Tokoh Wali Songo yang mengombinasikan misi dakwah dengan perluasan wilayah, penguatan armada militer maritim, diplomasi internasional, serta memegang kendali pemerintahan struktural sebagai Sultan Cirebon adalah \dots
    \begin{enumerate}[label=\alph*.]
        \item Sunan Giri
        \item Sunan Gunung Jati (Syarif Hidayatullah)
        \item Sunan Ampel
        \item Sunan Drajat
        \item Sunan Kudus
    \end{enumerate}

% ==========================================
% MATERI: Q.S. AN-NUR/24: 2 & PERGAULAN BEBAS
% ==========================================

\item Berdasarkan kandungan hukum Q.S. an-Nur/24: 2, sanksi syariat (hukuman dera/cambuk) yang dijatuhkan bagi pelaku zina perempuan dan laki-laki yang berstatus \textit{ghairu muhshan} (belum pernah terikat pernikahan yang sah) adalah sebanyak \dots
    \begin{enumerate}[label=\alph*.]
        \item 50 kali cambukan
        \item 75 kali cambukan
        \item 80 kali cambukan
        \item 100 kali cambukan
        \item 200 kali cambukan
    \end{enumerate}

\item Di era kontemporer modern, pergaulan bebas berkembang melalui berbagai manifestasi permisif yang melanggar batas moral. Bentuk aktivitas kedekatan fisik di luar pernikahan yang mengadopsi gaya hidup sekular-hedonis diidentifikasi sebagai \dots
    \begin{enumerate}[label=\alph*.]
        \item Hubungan pacaran tanpa batas moral
        \item Sistem ta'aruf pragmatis struktural
        \item Institusi pernikahan siri tanpa dokumentasi
        \item Perkawinan mut'ah kontemporer formal
        \item Hubungan silaturahmi komunal inklusif
    \end{enumerate}

\item Perbuatan zina dikategorikan sebagai perbuatan keji (\textit{fahisyah}). Dari dimensi kesehatan medis, perbuatan zina merupakan faktor pemicu utama penyebaran \dots
    \begin{enumerate}[label=\alph*.]
        \item Penyakit diabetes melitus kronis dan hipertensi sistemik
        \item Vektor utama Penyakit Menular Seksual (PMS) seperti HIV/AIDS dan sifilis
        \item Kerusakan degeneratif pada organ paru-paru akibat polusi udara
        \item Gangguan fungsi pencernaan akibat pola makan yang tidak teratur
        \item Penurunan imunitas tubuh yang disebabkan oleh faktor penuaan alami
    \end{enumerate}

% ==========================================
% MATERI: MAHABBATULLAH, KHAUF, RAJA', TAWAKAL
% ==========================================

\item Munculnya rasa kelezatan rohani, kedamaian hati (\textit{tumaninah}), serta kekhusyukan yang dalam saat menegakkan salat, berzikir, dan membaca Al-Qur'an merupakan indikasi dari \dots
    \begin{enumerate}[label=\alph*.]
        \item Tingginya tingkat kecerdasan intelektual kognitif manusia
        \item Tanda-tanda adanya rasa cinta sejati kepada Allah SWT (\textit{mahabbatullah})
        \item Karakter dasar manusia yang menyukai kesunyian lingkungan fisik
        \item Pengaruh eksternal dari kondisi sosiologis tempat tinggal individu
        \item Keberhasilan latihan meditasi pernapasan fisik yang teratur
    \end{enumerate}

\item Merenungkan keindahan, kompleksitas, dan keagungan penciptaan alam semesta sebagai cerminan sifat-sifat Allah yang Maha Agung disebut dengan metode \dots
    \begin{enumerate}[label=\alph*.]
        \item \textit{Tadakkur an-Ni'am}
        \item \textit{Tafakur fi al-Khalq}
        \item \textit{Mujahadatun Nafs}
        \item \textit{Amar Ma'ruf}
        \item \textit{Nahi Munkar}
    \end{enumerate}

\item Konsep \textit{khauf} (takut kepada Allah) yang benar dalam perspektif spiritual Islam dicirikan oleh perilaku hamba yang \dots
    \begin{enumerate}[label=\alph*.]
        \item Menjauhkan diri dari aktivitas peribadatan karena merasa tidak layak
        \item Berlari mendekat mencari perlindungan di balik tirai ampunan Allah dengan menjauhi maksiat
        \item Berdiam diri tanpa melakukan ikhtiar karena takut menghadapi kegagalan duniawi
        \item Membenci segala bentuk ketetapan takdir yang dirasa memberatkan kehidupan
        \item Menampilkan ketakutan fisik secara berlebihan di hadapan sesama manusia
    \end{enumerate}

\item Berikut ini yang merupakan tanda nyata di dalam diri seseorang yang memiliki sifat \textit{khauf} kepada kemurkaan Allah SWT adalah \dots
    \begin{enumerate}[label=\alph*.]
        \item Tercegahnya lisan dari perilaku ghibah dan dusta, serta terjaganya pandangan dari maksiat visual
        \item Suka menceritakan amal kebaikan diri sendiri agar dinilai sebagai hamba yang saleh
        \item Selalu menuntut keadilan hukum ditegakkan untuk orang lain tetapi abai pada diri sendiri
        \item Merasa aman dari ancaman siksa neraka karena tumpukan amal ibadah wajibnya
        \item Mengisolasi diri sepenuhnya dari dinamika interaksi sosial kemasyarakatan
    \end{enumerate}

\item Penyerahan total hasil akhir dari segala urusan kepada Allah SWT setelah melakukan usaha (ikhtiar) secara maksimal, cerdas, terencana, dan halal disebut dengan \dots
    \begin{enumerate}[label=\alph*.]
        \item \textit{Khauf}
        \item \textit{Raja'}
        \item \textit{Tawakkal}
        \item \textit{Syaja'ah}
        \item \textit{Mahabbah}
    \end{enumerate}

\item Karakteristik psikologis yang muncul sebagai manfaat utama dari penerapan sifat tawakal yang matang saat menghadapi kegagalan usaha adalah \dots
    \begin{enumerate}[label=\alph*.]
        \item Munculnya rasa putus asa yang mendalam terhadap rahmat Allah
        \item Ketenangan batin yang menghindarkan diri dari stres atau depresi akut
        \item Lahirnya sikap menyalahkan lingkungan eksternal sebagai penyebab utama
        \item Meningkatnya ambisi material untuk membalas kegagalan dengan segala cara
        \item Hilangnya motivasi untuk kembali melakukan perencanaan kerja di masa depan
    \end{enumerate}

\item Mengapa akal diposisikan sebagai elemen vital yang wajib dipelihara (\textit{hifzhu al-aql}) dalam sistem hukum Islam?
    \begin{enumerate}[label=\alph*.]
        \item Karena akal merupakan modal utama manusia dalam mengumpulkan kekayaan finansial
        \item Karena akal adalah basis pembebanan syariat (\textit{manathat taklif}) yang membedakan derajat manusia dengan makhluk lain
        \item Karena akal berfungsi untuk menggantikan peran wahyu ilahi dalam penentuan hukum fikih
        \item Karena akal tidak dapat dipengaruhi oleh bisikan emosional negatif maupun nafsu internal
        \item Karena akal merupakan satu-satunya alat penentu keselamatan mutlak manusia di akhirat
    \end{enumerate}

% ==========================================
% TAMBAHAN (SOAL KE-41 SD 60 SESUAI KISI-KISI)
% ==========================================

\item Berdasarkan petunjuk sabda Rasulullah Saw., salah satu metode pertolongan pertama secara fisik dan spiritual untuk meredam sifat temperamental (\textit{ghadhab}) yang sedang memuncak adalah \dots
    \begin{enumerate}[label=\alph*.]
        \item Berbicara dengan suara yang lebih keras untuk melampiaskan emosi
        \item Segera mengubah posisi tubuh (jika berdiri maka duduk) dan mengambil air wudu
        \item Meninggalkan tempat kejadian dan memutus tali silaturahmi
        \item Melakukan aktivitas fisik yang berat agar lelah
        \item Membalas perkataan lawan bicara agar mendapatkan kepuasan batin
    \end{enumerate}

\item Sikap \textit{syaja'ah} (berani) sangat erat kaitannya dengan kejujuran (\textit{ash-shidqu}). Contoh nyata keterkaitan antara \textit{syaja'ah} dan kejujuran di lingkungan sekolah adalah \dots
    \begin{enumerate}[label=\alph*.]
        \item Berani mengakui kesalahan diri sendiri dan tidak menyontek saat ujian meskipun tidak diawasi
        \item Berani membantah teguran guru demi membela teman satu kelompok
        \item Berani menantang perkelahian pelajar dari sekolah lain demi menjaga harga diri
        \item Berani berbohong demi melindungi rahasia buruk sahabat dekat
        \item Berani melanggar tata tertib sekolah sebagai bentuk kebebasan berekspresi
    \end{enumerate}

\item \textit{Mujahadatun nafs} (kontrol diri) mengharuskan seorang muslim untuk senantiasa melakukan \textit{muhasabah}. Arti dari \textit{muhasabah} dalam konteks pengendalian diri adalah \dots
    \begin{enumerate}[label=\alph*.]
        \item Merencanakan masa depan dengan ambisi yang tinggi
        \item Mengingat-ingat kesalahan orang lain untuk dijadikan pelajaran
        \item Mengevaluasi dan mengintrospeksi niat serta perbuatan diri sendiri setiap saat
        \item Menahan rasa lapar dan haus secara berkesinambungan tanpa henti
        \item Memisahkan diri dari interaksi sosial masyarakat (uzlah) secara permanen
    \end{enumerate}

\item Urutan hierarki yang paling tepat dan disepakati oleh mayoritas ulama Ushul Fikih terkait prioritas perlindungan dalam \textit{Al-Kulliyatu al-Khamsah} (lima prinsip dasar) adalah \dots
    \begin{enumerate}[label=\alph*.]
        \item Harta, Akal, Keturunan, Jiwa, Agama
        \item Agama, Jiwa, Akal, Keturunan, Harta
        \item Jiwa, Agama, Akal, Harta, Keturunan
        \item Keturunan, Harta, Jiwa, Agama, Akal
        \item Agama, Akal, Jiwa, Harta, Keturunan
    \end{enumerate}

\item Di era pergaulan digital, tantangan terbesar dalam memelihara agama (\textit{hifzhu al-din}) adalah masifnya informasi yang menyesatkan. Tindakan paling tepat bagi seorang pelajar muslim untuk menjaga agamanya adalah \dots
    \begin{enumerate}[label=\alph*.]
        \item Menerima semua informasi keagamaan dari internet tanpa perlu verifikasi
        \item Mempelajari ilmu agama (sanad keilmuan) melalui guru atau ulama yang tepercaya
        \item Berhenti menggunakan media sosial secara total
        \item Menyebarkan informasi apa saja asalkan memiliki unsur bahasa Arab
        \item Berdebat secara agresif dengan siapapun yang berbeda pandangan di media sosial
    \end{enumerate}

\item Tindakan mengakhiri hidup sendiri (bunuh diri) dengan alasan putus asa menghadapi cobaan hidup sangat diharamkan dalam Islam. Pengharaman ini merupakan penegasan dari prinsip \dots
    \begin{enumerate}[label=\alph*.]
        \item \textit{Hifzhu al-Din}
        \item \textit{Hifzhu al-Aql}
        \item \textit{Hifzhu al-Nafs}
        \item \textit{Hifzhu al-Nasl}
        \item \textit{Hifzhu al-Maal}
    \end{enumerate}

\item Syariat Islam sangat melarang keras perbuatan aborsi (pengguguran kandungan) tanpa adanya indikasi kedaruratan medis yang mengancam nyawa ibu. Larangan tersebut merupakan bentuk implementasi dari \dots
    \begin{enumerate}[label=\alph*.]
        \item \textit{Hifzhu al-Nasl} (Menjaga Keturunan) dan \textit{Hifzhu al-Nafs} (Menjaga Jiwa)
        \item \textit{Hifzhu al-Aql} (Menjaga Akal) dan \textit{Hifzhu al-Maal} (Menjaga Harta)
        \item \textit{Hifzhu al-Din} (Menjaga Agama) dan \textit{Hifzhu al-Aql} (Menjaga Akal)
        \item \textit{Hifzhu al-Maal} (Menjaga Harta) dan \textit{Hifzhu al-Din} (Menjaga Agama)
        \item \textit{Hifzhu al-Irdh} (Menjaga Kehormatan) semata
    \end{enumerate}

\item Salah satu bentuk \textit{hifzhu al-maal} (menjaga harta) dari segi pengembangannya (\textit{min janib al-wujud}) adalah \dots
    \begin{enumerate}[label=\alph*.]
        \item Menjatuhkan hukuman potong tangan bagi pencuri
        \item Mewajibkan zakat dan menganjurkan sedekah, infak, serta perniagaan yang halal
        \item Mengharamkan praktik riba dalam setiap transaksi perbankan
        \item Melarang tindak pidana korupsi dan suap (risywah)
        \item Mencegah terjadinya penipuan (gharar) dalam jual beli daring
    \end{enumerate}

\item Salah satu rahasia kesuksesan dakwah Wali Songo di Nusantara adalah menjunjung tinggi nilai \textit{tasamuh}. Makna dari \textit{tasamuh} dalam metode dakwah tersebut adalah \dots
    \begin{enumerate}[label=\alph*.]
        \item Sikap keras kepala dalam berdebat
        \item Sikap toleransi, menghargai, dan moderat terhadap keragaman budaya setempat
        \item Memaksa masyarakat untuk meninggalkan kebudayaan asli mereka secara instan
        \item Mengasingkan diri dari hiruk-pikuk kehidupan duniawi
        \item Menguasai jalur ekonomi untuk menekan masyarakat bawah
    \end{enumerate}

\item Sunan Kalijaga dikenal menggunakan media kesenian lokal untuk memasukkan nilai-nilai Islam. Istilah "Jimat Kalimasada" dalam lakon pewayangan gubahan beliau sebenarnya merepresentasikan \dots
    \begin{enumerate}[label=\alph*.]
        \item Kekuatan magis dari senjata tradisional Jawa
        \item Dua Kalimat Syahadat (\textit{Kalimatain Syahadatain})
        \item Lima rukun Islam yang diwajibkan
        \item Rapalan mantra untuk memanggil hujan
        \item Silsilah keturunan raja-raja Majapahit
    \end{enumerate}

\item Ajaran falsafah \textit{"Moh Limo"} yang diajarkan oleh Sunan Ampel berfokus pada perbaikan moral masyarakat. Kata \textit{"Moh Maling"} memiliki arti tidak mau \dots
    \begin{enumerate}[label=\alph*.]
        \item Mabuk-mabukan
        \item Bermain judi
        \item Mencuri barang orang lain
        \item Menghisap candu/narkoba
        \item Melakukan perzinaan
    \end{enumerate}

\item Pusat pergerakan dakwah Sunan Gunung Jati difokuskan di wilayah barat pulau Jawa. Berkat strategi dakwah dan perannya, beliau berhasil meletakkan dasar kekuatan politik Islam di wilayah \dots
    \begin{enumerate}[label=\alph*.]
        \item Demak dan Kudus
        \item Gresik dan Tuban
        \item Banten dan Cirebon
        \item Ampeldenta dan Lamongan
        \item Giri Kedaton dan Madura
    \end{enumerate}

\item Dalam Q.S. An-Nur/24: 2, terdapat kalimat \textit{"wala ta'khudzkum bihima ra'fatun fi dinillahi"}. Ayat ini memerintahkan kepada penegak hukum untuk \dots
    \begin{enumerate}[label=\alph*.]
        \item Memberikan keringanan hukuman jika pelaku berjanji tidak mengulangi
        \item Tidak merasa belas kasihan dalam menjalankan hukuman syariat (hukum Allah) kepada pezina
        \item Menutupi aib pezina jika ia berasal dari keluarga yang terhormat
        \item Membebaskan pelaku zina jika keduanya melakukan atas dasar suka sama suka
        \item Menyerahkan urusan hukuman sepenuhnya kepada masyarakat secara main hakim sendiri
    \end{enumerate}

\item Perzinaan tidak terjadi secara tiba-tiba, melainkan diawali oleh pintu-pintu pergaulan bebas. Oleh karena itu, Islam mengajarkan langkah preventif (\textit{sad az-zari'ah}), yang di antaranya adalah larangan \textit{khalwat}. Pengertian \textit{khalwat} adalah \dots
    \begin{enumerate}[label=\alph*.]
        \item Berdua-duaan antara laki-laki dan perempuan yang bukan mahram di tempat sepi
        \item Membuka aurat di depan umum
        \item Memandang lawan jenis dengan syahwat
        \item Mengonsumsi minuman keras bersama-sama
        \item Berkomunikasi melalui media sosial
    \end{enumerate}

\item Dampak destruktif berskala besar dari maraknya perzinaan di tengah masyarakat dari sudut pandang sosial dan \textit{hifzhu al-nasl} (nasab) adalah \dots
    \begin{enumerate}[label=\alph*.]
        \item Peningkatan laju pertumbuhan ekonomi makro
        \item Tercampurnya nasab keturunan yang berakibat kacaunya hukum waris dan perwalian nikah
        \item Terciptanya generasi yang memiliki kreativitas intelektual tinggi
        \item Melemahnya stabilitas politik luar negeri sebuah negara
        \item Berkurangnya jumlah tindak kriminal pencurian dan perampokan
    \end{enumerate}

\item Tingkatan tertinggi dalam pelaksanaan ibadah adalah ketika seorang hamba menyembah Allah Swt. bukan semata-mata karena mengharap surga atau takut neraka, melainkan didasari oleh rasa cinta (\textit{mahabbatullah}). Karakteristik ibadah ini disebut tingkatan \dots
    \begin{enumerate}[label=\alph*.]
        \item \textit{Ibadatul 'Ubbad} (Ibadahnya hamba sahaya)
        \item \textit{Ibadatut Tujjar} (Ibadahnya pedagang)
        \item \textit{Ibadatul Ahrar} (Ibadahnya orang merdeka/pecinta)
        \item \textit{Ibadatul Kharijin} (Ibadahnya kaum pemberontak)
        \item \textit{Ibadatul Ghafilin} (Ibadahnya orang yang lalai)
    \end{enumerate}

\item Rasa takut kepada Allah (\textit{khauf}) harus seimbang dengan rasa harap (\textit{raja'}). Jika seseorang hanya memiliki sifat \textit{khauf} tanpa \textit{raja'}, maka bahaya psikologis yang akan menimpanya adalah \dots
    \begin{enumerate}[label=\alph*.]
        \item Bersikap sombong karena merasa paling suci
        \item Terjerumus ke dalam sifat putus asa (\textit{qunut}) dari rahmat Allah
        \item Merasa aman dari maksiat karena Allah Maha Pengampun
        \item Menjadi pribadi yang terlalu toleran terhadap dosa
        \item Memiliki optimisme yang buta tanpa amal nyata
    \end{enumerate}

\item Perhatikan ilustrasi berikut: "Andi belajar dengan tekun setiap hari, berdoa di sepertiga malam, lalu saat ujian ia mengerjakannya dengan jujur. Setelah ujian selesai, ia menyerahkan hasil kelulusannya kepada Allah." Sikap Andi merupakan bentuk pengamalan dari \dots
    \begin{enumerate}[label=\alph*.]
        \item \textit{Jabariyah} (Pasrah tanpa usaha)
        \item \textit{Qadariyah} (Usaha tanpa campur tangan Tuhan)
        \item \textit{Tawakal al-Haq} (Tawakal yang benar)
        \item \textit{Tawakal as-Su'} (Tawakal yang keliru)
        \item \textit{Tasahul} (Meremehkan urusan)
    \end{enumerate}

\item Salah satu ciri mendasar orang yang bertawakal dengan benar adalah ia tidak akan merasa \dots ketika impian dan target pekerjaannya belum tercapai sesuai rencana.
    \begin{enumerate}[label=\alph*.]
        \item Termotivasi
        \item Syukur
        \item Sabar
        \item Kecewa dan frustrasi berlebihan
        \item Semangat untuk mengevaluasi diri
    \end{enumerate}

\item Perkembangan literasi digital mewajibkan kita untuk menyaring informasi dan menjauhi berita hoaks. Sikap kritis (\textit{tabayyun}) dalam menerima berita merupakan bentuk perlindungan kontemporer terhadap prinsip \dots
    \begin{enumerate}[label=\alph*.]
        \item \textit{Hifzhu al-Nafs}
        \item \textit{Hifzhu al-Maal}
        \item \textit{Hifzhu al-Nasl}
        \item \textit{Hifzhu al-Aql}
        \item \textit{Hifzhu al-Din}
    \end{enumerate}

\end{enumerate}
\end{document}